\begin{table}[htbp]
\centering
\caption{Illustrative comparison of methodological coverage in related work}
\label{tab:relatedwork_gap_template}
\setlength{\tabcolsep}{4pt}
renewcommand{\arraystretch}{1.2}

\begin{threeparttable}
\begin{tabularx}{\textwidth}{p{2.4cm} X *{5}{>{\centering\arraybackslash}p{1.5cm}}}
\toprule
\textbf{Category} &
\textbf{Representative Systems} &
\textbf{Documentation support} &
\textbf{Communication guidance} &
\textbf{Domain-specific strategies} &
\textbf{Active consultation guidance} &
\textbf{Source grounding} \\
\midrule
AI Documentation & Doctolib, Suki.ai & \yes & \no & \no & \no & \no \\
\midrule
Clinical Decision Support & UpToDate, OpenEvidence & \no & \no & \partialsymbol & \no & \partialsymbol \\
\midrule
LLM Medical Assistants & Glass Health & \yes & \partialsymbol & \no & \no & \no \\
\midrule
Communication Frameworks & Ariadne Labs Toolkit & \no & \yes & \yes & \partialsymbol & \yes \\
\midrule
\textbf{This project} & \lightexample{(Navigation Assistant)} & \yes & \yes & \yes & \yes & \yes \\
\bottomrule
\end{tabularx}

\begin{tablenotes}\small
\item \yes\ indicates full coverage of the property.
\item \partialsymbol\ indicates partial or limited coverage.
\item \no\ indicates absence of evidence in the cited work.
\item Source grounding refers to explicit citation and attribution of evidence-based resources.
\end{tablenotes}
\end{threeparttable}
\end{table}
